% !TEX root = ../main.tex
% CHU SO HUU: Ngọc Trâm | KHONG THEO PREAMBLE, KHONG BEGIN DOCUMENT TRONG THAN PHAN NAY.
\documentclass[../main.tex]{subfiles}
\begin{document}
\section{Kết quả kiểm định hai mẫu độc lập}
\label{sec:t35:result}

Phân tích sử dụng toàn bộ 125\,973 bản ghi KDDTrain+ hợp lệ, gồm 67\,343 Normal và
58\,630 Attack. Ba đặc trưng được kiểm định trên thang gốc bằng Welch hai phía tại
$\alpha=0{,}05$; ba p-value chính được hiệu chỉnh Holm. Không có bản ghi nào bị loại
do thiếu, NaN, vô hạn hoặc sai miền. Các p-value bị underflow số học được trình bày
là $p<10^{-300}$, không viết $p=0$.

\subsection{Thống kê mô tả}
\label{sec:t35:descriptive}

Bảng~\ref{tab:t35:descriptive} cho thấy hai nhóm chỉ mất cân bằng nhẹ về cỡ
mẫu (tỷ số 1,149), nhưng vị trí và độ phân tán khác nhau rõ. Đặc biệt, phương sai của
Attack lớn hơn Normal ở cả ba biến.

Ở \texttt{srv\_serror\_rate}, trung vị chuyển từ 0 ở Normal sang 1 ở Attack và tỷ lệ
giá trị bằng 1 của Attack đạt 58,920\% (Normal: 0,489\%). Với
\texttt{srv\_rerror\_rate}, cả hai trung vị đều bằng 0, nhưng Attack có tỷ lệ bằng 1
là 20,607\%, cao hơn Normal (4,076\%). Vì vậy, chênh lệch trung bình của RQ2 không
đồng nghĩa mọi quan sát Attack đều có lỗi REJ. Đối với \texttt{count}, khoảng tứ phân
vị Attack (40--241) dịch lên rõ so với Normal (1--14), nhưng hai nhóm vẫn cùng miền
0--511.

Không có nhãn tấn công nào bị loại khỏi Attack.

\subsection{Phân tích phân phối và Q--Q plot}
\label{sec:t35:qq}

Các Q--Q plot dùng mẫu cố định tối đa 1\,000 quan sát/nhóm để tránh chồng hàng chục
nghìn điểm; đây chỉ là bước vẽ, không thay thế dữ liệu kiểm định.

\begin{figure}[H]
\centering
\includegraphics[width=0.97\linewidth]{figures/task35_qq_srv_serror_rate.png}
\caption{Q--Q plot của \texttt{srv\_serror\_rate}. Các dải ngang tại 0 và 1 phản ánh
khối lượng xác suất tại hai biên; mỗi panel dùng tối đa 1\,000 quan sát.}
\label{fig:t35:qq-serror}
\end{figure}

Hình~\ref{fig:t35:qq-serror} không tạo thành đám mây gần đường thẳng. Ở Normal,
phần lớn điểm nằm trên dải $y=0$, sau đó tăng đột ngột ở đuôi phải; skewness đạt 9,169.
Ở Attack, hai dải dày tại 0 và 1 tạo dạng gần nhị điểm, phù hợp với tỷ lệ bằng 0 là
40,486\% và bằng 1 là 58,920\%. Hai hình dạng khác nhau rõ và đều không phù hợp với
phân phối chuẩn của quan sát gốc.

\begin{figure}[H]
\centering
\includegraphics[width=0.97\linewidth]{figures/task35_qq_srv_rerror_rate.png}
\caption{Q--Q plot của \texttt{srv\_rerror\_rate}. Các bậc và dải ngang xuất hiện do
biến tỷ lệ rời rạc, tập trung tại 0 và 1.}
\label{fig:t35:qq-rerror}
\end{figure}

Với \texttt{srv\_rerror\_rate}, Hình~\ref{fig:t35:qq-rerror} cũng có dải dài
tại 0, nhưng phần giá trị dương của Normal cong mạnh và nhảy lên dải 1 ở đuôi phải
(skewness 4,426). Attack có nhiều bậc trung gian hơn trước khi đạt 1 và vẫn lệch phải
(skewness 1,434). Đây không phải bản sao hình dạng RQ1: ở RQ2, trung vị và Q3 đều bằng
0 trong cả hai nhóm, cho thấy khối tại 0 chi phối mạnh hơn phần trung tâm.

\begin{figure}[H]
\centering
\includegraphics[width=0.97\linewidth]{figures/task35_qq_count.png}
\caption{Q--Q plot của \texttt{count}. Normal có đuôi phải dài; Attack có dạng bậc
và trần tại 511, thể hiện tính rời rạc và bị chặn của biến đếm.}
\label{fig:t35:qq-count}
\end{figure}

Hình~\ref{fig:t35:qq-count} cho thấy Normal cong lên rất mạnh ở nửa phải và có
đuôi tới trên 400 trong mẫu vẽ, phù hợp skewness 3,878. Attack gần tuyến tính hơn ở
phần giữa nhưng lệch khỏi đường tham chiếu ở cả hai đuôi, có các bậc và dải ngang tại
511; skewness vẫn là 0,747. Q--Q plot vì vậy bác bỏ mô tả chuẩn đơn giản cho dữ liệu
gốc, đồng thời cho thấy dạng sai lệch của hai nhóm không giống nhau.

Q--Q plot chỉ là chẩn đoán trực quan. Cỡ mẫu lớn hỗ trợ xấp xỉ của trung bình mẫu,
nhưng không xóa khối biên, tính rời rạc hay phụ thuộc tiềm ẩn. Bảng
\ref{tab:t35:assumptions} bổ sung Brown--Forsythe và tỷ số phương sai; tất cả
đều cho thấy Student t-test phương sai gộp không thích hợp.

\begin{table}[H]
\centering
\caption{Chẩn đoán phương sai và hình dạng phân phối}
\label{tab:t35:assumptions}
\small
\setlength{\tabcolsep}{4pt}
\begin{tabularx}{\textwidth}{lrrrX}
\toprule
Đặc trưng & Brown--Forsythe & p-value & Tỷ số phương sai & Chẩn đoán trọng tâm \\
\midrule
\texttt{srv\_serror\_rate} & 42\,179,95 & $<10^{-300}$ & 32,175 & Bị chặn; dồn mạnh tại 0 và 1. \\
\texttt{srv\_rerror\_rate} & 8\,651,41 & $<10^{-300}$ & 3,999 & Q1=Q3=0; lệch phải và nhiều giá trị trùng. \\
\texttt{count} & 46\,458,65 & $<10^{-300}$ & 5,296 & Biến đếm; Normal đuôi phải dài, hai nhóm cùng trần 511. \\
\bottomrule
\end{tabularx}
\end{table}

\subsection{Kết quả kiểm định tỷ lệ lỗi SYN}
\label{sec:t35:rq1}

RQ1 xem xét \texttt{srv\_serror\_rate}. Bằng chứng mô tả cho thấy Normal có trung
bình 0,01208 (KTC trung bình nhóm 95\%: 0,01143--0,01274), còn Attack là 0,59307
(0,58910--0,59704). Phương sai Attack lớn gấp 32,175 lần Normal và Q--Q plot có khối
biên rõ, nên Welch được giữ làm kiểm định chính thay vì Student t-test phương sai gộp.

Khoảng tin cậy hoàn toàn âm và p-value Holm nhỏ hơn $10^{-300}$, do đó bác bỏ
$H_0$. Theo chiều định nghĩa Normal trừ Attack, kết quả tương ứng Attack cao hơn khoảng
0,581 đơn vị tỷ lệ. Hedges $g=-1{,}7070$ (KTC xấp xỉ
$[-1{,}7199;-1{,}6940]$) cho thấy chênh lệch chuẩn hóa đáng kể trong mẫu. Mann--Whitney
và Yuen--Welch 20\% đều cho $p<10^{-300}$, củng cố độ bền của kết luận về khác biệt,
nhưng không biến Welch thành mô hình hoàn chỉnh cho biến tỷ lệ dồn biên. Đặc trưng phản
ánh lỗi bắt tay TCP theo dịch vụ, không chứng minh lỗi SYN là nguyên nhân của tấn công.

\subsection{Kết quả kiểm định tỷ lệ lỗi REJ}
\label{sec:t35:rq2}

RQ2 phân tích \texttt{srv\_rerror\_rate}. Normal có trung bình 0,04463, Attack là
0,20911; dù vậy, trung vị và Q3 của cả hai nhóm đều bằng 0. Sự kết hợp giữa khối tại 0
và phương sai Attack gấp 3,999 lần Normal là lý do không diễn giải chênh lệch trung
bình như khác biệt đồng đều ở mọi kết nối.

Bác bỏ $H_0$ sau hiệu chỉnh Holm. Attack cao hơn Normal khoảng 0,1645 đơn vị tỷ lệ;
Hedges $g=-0{,}5254$ (KTC xấp xỉ $[-0{,}5366;-0{,}5141]$) nhỏ hơn về trị tuyệt đối
so với RQ1. Mann--Whitney cho $p<10^{-300}$; Yuen--Welch cho
$p=4{,}905\times10^{-8}$ với chênh lệch trung bình cắt ngọn chỉ $-0{,}01519$. Sự
giảm mạnh sau cắt ngọn cho thấy chênh lệch trung bình gốc chịu ảnh hưởng đáng kể của
khối giá trị lớn, dù bằng chứng về khác biệt vẫn còn. Đây là tín hiệu về lỗi REJ theo
dịch vụ, không phải bằng chứng rằng một ngưỡng đơn biến đủ nhận diện Attack.

\subsection{Kết quả kiểm định số lượng kết nối gần đây}
\label{sec:t35:rq3}

RQ3 xét \texttt{count}. Normal có trung bình 22,518 và trung vị 4, trong khi Attack
có trung bình 154,850 và trung vị 138. Chênh lệch của cả trung bình lẫn trung vị cùng
hướng, nhưng độ lệch phải, ngoại lệ và trần 511 cho thấy phân bố không chuẩn.

Khoảng tin cậy cho thấy Attack cao hơn Normal khoảng 131,246--133,418 kết nối trung
bình trong định nghĩa cửa sổ của NSL-KDD. Bác bỏ $H_0$ sau Holm; Hedges
$g=-1{,}4143$ (KTC xấp xỉ $[-1{,}4266;-1{,}4019]$). Mann--Whitney và Yuen--Welch
đều có $p<10^{-300}$; chênh lệch trung bình cắt ngọn là $-140{,}988$. Kết quả phù hợp
với việc nhiều hành vi Attack tạo kết nối tập trung trong thời gian ngắn, nhưng nhóm
Attack là hỗn hợp bốn họ và hai miền giá trị vẫn chồng lấp; không thể suy ra mọi dạng
tấn công đều có \texttt{count} cao.

\subsection{So sánh tổng hợp ba kiểm định}
\label{sec:t35:summary}

\begin{figure}[H]
\centering
\includegraphics[width=0.98\linewidth]{figures/task35_mean_difference_forest.png}
\caption{Chênh lệch trung bình Normal $-$ Attack và KTC 95\% theo Welch. Ba panel
dùng trục riêng vì đặc trưng có đơn vị khác nhau; đường đứt đoạn tại 0 là mốc không
chênh lệch.}
\label{fig:t35:mean-difference}
\end{figure}

\begin{figure}[H]
\centering
\includegraphics[width=0.80\linewidth]{figures/task35_effect_size_forest.png}
\caption{Hedges $g$ và KTC 95\% xấp xỉ cho ba RQ. Trị tuyệt đối lớn không được diễn
giải là năng lực dự đoán ngoài mẫu.}
\label{fig:t35:effect-size}
\end{figure}

Cả ba giả thuyết bị bác bỏ sau Holm, nhưng mức độ và hình dạng khác biệt không giống
nhau. RQ1 và RQ3 có chênh lệch chuẩn hóa lớn trong mẫu; RQ2 nhỏ hơn và nhạy hơn với
cắt ngọn. Các biến tỷ lệ gây khó khăn vì khối biên, còn \texttt{count} gây khó khăn
do tính rời rạc, lệch phải và trần dữ liệu. Do cỡ mẫu rất lớn, p-value cực nhỏ phải
được đọc cùng khoảng tin cậy, effect size và phân bố chồng lấp.

Kết quả bổ sung cho Pearson: hệ số Pearson hiện có giữa
\path{srv_serror_rate} và \path{dst_host_srv_serror_rate} là 0,9863; giữa
\path{srv_rerror_rate} và \path{dst_host_srv_rerror_rate} là 0,9702.
Các hệ số này cho biết tính dư thừa tuyến tính giữa đặc trưng, còn Welch cho biết
chênh lệch trung bình theo nhãn. ANOVA năm nhóm đối với \texttt{count} trả lời câu hỏi
khác: nó giữ riêng DoS, Probe, R2L, U2R, trong khi phân tích này gộp chúng thành Attack.
Hai phân tích dùng cùng bản ghi nên không độc lập về thông tin.

\FloatBarrier

\end{document}
